% Standardized TikZ commands for Pneumatic Elements
\newcommand{\cySingle}{
    \draw (-2, 0) rectangle (2, 1);
    \draw (-1, 0) -- (-1, 1); \draw (-0.7, 0) -- (-0.7, 1);
    \draw (-0.7, 0.4) -- (2.8, 0.4); \draw (-0.7, 0.6) -- (2.8, 0.6);
    \draw (-1.5, 0) -- (-1.5, -0.5);
    % Spring
    \draw (-0.7, 0.2) -- (-0.2, 0.8) -- (0.3, 0.2) -- (0.8, 0.8) -- (1.3, 0.2) -- (1.8, 0.8) -- (2, 0.5);
}

\newcommand{\cyDouble}{
    \draw (-2, 0) rectangle (2, 1);
    \draw (-0.5, 0) -- (-0.5, 1); \draw (-0.2, 0) -- (-0.2, 1);
    \draw (-0.2, 0.4) -- (2.8, 0.4); \draw (-0.2, 0.6) -- (2.8, 0.6);
    \draw (-1, 0) -- (-1, -0.5); \draw (1.5, 0) -- (1.5, -0.5);
}

\newcommand{\valveThreeTwo}{
    \draw (0,0) rectangle (1,1); \draw (1,0) rectangle (2,1);
    % Left box (actuated)
    \draw[->] (0.2, 0) -- (0.2, 1);
    \draw (0.8, 0) -- (0.8, 0.2); \draw (0.65, 0.2) -- (0.95, 0.2);
    % Right box (normal)
    \draw (1.2, 0) -- (1.2, 0.2); \draw (1.05, 0.2) -- (1.35, 0.2);
    \draw[->] (1.2, 1) -- (1.8, 0); 
}

\newcommand{\portSupply}[2]{
    \begin{scope}[shift={(#1, #2)}]
        \draw (0, 0) -- (0, -0.5);
        \draw (-0.08, -0.8) -- (0.08, -0.8) -- (0, -0.5) -- cycle;
    \end{scope}
}

\newcommand{\portExhaust}[2]{
    \begin{scope}[shift={(#1, #2)}]
        \draw (0, 0) -- (0, -0.5);
        \draw (-0.08, -0.5) -- (0.08, -0.5) -- (0, -0.8) -- cycle;
    \end{scope}
}

\newcommand{\valveFiveTwo}{
    \draw (0,0) rectangle (1,1); \draw (1,0) rectangle (2,1);
    
    % Left box (actuated)
    \draw (0.3, 0) -- (0.3, 0.2); \draw (0.2, 0.2) -- (0.4, 0.2);
    \draw[->] (0.5, 0) -- (0.3, 1);
    \draw[->] (0.7, 1) -- (0.7, 0);
    
    % Right box (normal)
    \draw[->] (1.3, 1) -- (1.3, 0);
    \draw[->] (1.5, 0) -- (1.7, 1);
    \draw (1.7, 0) -- (1.7, 0.2); \draw (1.6, 0.2) -- (1.8, 0.2);
}

\newcommand{\actButton}{
    \draw (0, 0.8) -- (-0.7, 0.8);
    \draw (0, 0.2) -- (-0.7, 0.2);
    \draw (-0.7, 0) -- (-0.7, 1.0);
}

\newcommand{\actRoller}{
    \draw (0, 0.5) -- (-0.3, 0.5); \draw (-0.3, 0.3) circle (0.2);
}

\newcommand{\actCam}{
    \actRoller
    \draw (-0.5, 0.1) -- (-0.5, 0.9);
    \draw (-0.5, 0.9) -- (-0.8, 0.6);
    \draw (-0.5, 0.7) -- (-0.8, 0.4);
    \draw (-0.5, 0.5) -- (-0.8, 0.2);
    \draw (-0.5, 0.3) -- (-0.8, 0.0);
}

\newcommand{\sensorCapacitive}{
    \draw (-0.5, -0.5) rectangle (1.2, 0.5);
    \draw (0, 0.3) -- (0.15, 0.15) -- (0, 0) -- (-0.15, 0.15) -- cycle;
    \draw (0.4, 0.25) -- (0.4, 0.05); \draw (0.6, 0.25) -- (0.6, 0.05);
    \draw (0.2, 0.15) -- (0.4, 0.15); \draw (0.8, 0.15) -- (0.6, 0.15);
    \draw (0.9, 0.1) -- (0.9, 0.3) -- (1.1, 0.3);
}

\newcommand{\relayCoilDelay}{
    \draw (-0.5, -0.4) rectangle (0.5, 0.4);
    \draw (0, -0.4) -- (0, 0.4);
    \draw (-0.5, -0.4) -- (0, 0.4);
    \draw (-0.5, 0.4) -- (0, -0.4);
    \draw (0, 1.0) -- (0, 0.4);
    \draw (0, -0.4) -- (0, -1.0);
}

\newcommand{\actSpring}{
    \draw (2, 0.5) -- (2.1, 0.5);
    \draw[decorate, decoration={zigzag, segment length=2pt, amplitude=1.5pt}] (2.1, 0.5) -- (2.5, 0.5);
}

\newcommand{\actPilotLeft}{
    \draw (0, 0.5) -- (-0.3, 0.5);
    \draw[fill=white] (-0.5, 0.3) -- (-0.5, 0.7) -- (-0.3, 0.5) -- cycle;
}

\newcommand{\actPilotRight}{
    \draw (2, 0.5) -- (2.3, 0.5);
    \draw[fill=white] (2.5, 0.3) -- (2.5, 0.7) -- (2.3, 0.5) -- cycle;
}

\newcommand{\actSolenoidLeft}[1]{
    \draw (0, 0.5) -- (-0.5, 0.5); 
    \draw (-1.0, 0.3) rectangle (-0.5, 0.7); \draw (-1.0, 0.3) -- (-0.5, 0.7);
    \node[below] at (-0.75, 0.3) {#1};
}

\newcommand{\valveAnd}{
    \draw (-1, -0.5) rectangle (1, 0.5);
    \draw (-0.5, -0.5) -- (-0.5, -0.2); \draw (0.5, -0.5) -- (0.5, -0.2);
    \draw (0, 0.5) -- (0, 0.2);
    \draw (-0.8, -0.2) -- (0.8, -0.2);
    \draw (-0.8, -0.2) -- (-0.8, 0.2);
    \draw (0.8, -0.2) -- (0.8, 0.2);
    \draw (-0.8, 0.2) -- (0.8, 0.2);
    \node at (0, 0) {\&};
}

% Electrical contact: NO (normally open relay contact)
% Usage: \contactNO{x}{ytop}{label}
\newcommand{\contactNO}[3]{
    % NO Contact slash: from ytop-0.9 to ytop
    \draw (#1, {#2-0.9}) -- ({#1+0.4}, {#2});
    % Label: dashed line touches slash at center-y
    \draw[dashed] ({#1-0.4}, {#2-0.45}) -- ({#1+0.2}, {#2-0.45});
    \node[anchor=east] at ({#1-0.5}, {#2-0.45}) {#3};
}

% Electrical contact: NC (normally closed relay contact)
% Usage: \contactNC{x}{ytop}{label}
\newcommand{\contactNC}[3]{
    % NC Contact: Horizontal bar and crossing diagonal slash
    \draw ({#1-0.3}, {#2}) -- ({#1+0.6}, {#2});
    \draw (#1, {#2-0.9}) -- ({#1+0.6}, {#2+0.1});
    % Label
    \draw[dashed] ({#1-0.4}, {#2-0.45}) -- ({#1+0.3}, {#2-0.45});
    \node[anchor=east] at ({#1-0.5}, {#2-0.45}) {#3};
}

% Electrical switch: NO (normally open pushbutton)
\newcommand{\switchNO}[3]{
    \draw (#1, {#2-0.9}) -- ({#1+0.4}, {#2});
    % Pushbutton elements
    \draw[dashed] ({#1-0.4}, {#2-0.45}) -- ({#1+0.2}, {#2-0.45});
    \draw ({#1-0.6}, {#2-0.2}) -- ({#1-0.4}, {#2-0.2});
    \draw ({#1-0.6}, {#2-0.7}) -- ({#1-0.4}, {#2-0.7});
    \draw ({#1-0.6}, {#2-0.2}) -- ({#1-0.6}, {#2-0.7});
    \node[anchor=east] at ({#1-0.7}, {#2-0.45}) {#3};
}

% Electrical switch: NC (normally closed pushbutton)
\newcommand{\switchNC}[3]{
    \draw ({#1-0.3}, {#2}) -- ({#1+0.6}, {#2});
    \draw (#1, {#2-0.9}) -- ({#1+0.6}, {#2+0.1});
    % Pushbutton elements
    \draw[dashed] ({#1-0.4}, {#2-0.45}) -- ({#1+0.3}, {#2-0.45});
    \draw ({#1-0.6}, {#2-0.2}) -- ({#1-0.4}, {#2-0.2});
    \draw ({#1-0.6}, {#2-0.7}) -- ({#1-0.4}, {#2-0.7});
    \draw ({#1-0.6}, {#2-0.2}) -- ({#1-0.6}, {#2-0.7});
    \node[anchor=east] at ({#1-0.7}, {#2-0.45}) {#3};
}
% Relay coil
\newcommand{\relayCoil}[3]{
    \draw ({#1-0.5}, {#2}) rectangle ({#1+0.5}, {#2-1.0});
    \node at (#1, {#2-0.5}) {#3};
    % Add small dash for label connection like FluidSIM if needed, 
    % but standard rectangle with label inside is usually preferred.
}

% Solenoid coil (FluidSIM style: slashed rectangle + hourglass valve symbol)
% Usage: \solenoidCoil{x}{ytop}{label}
\newcommand{\solenoidCoil}[3]{
    % Slashed rectangle (1 unit tall)
    \draw ({#1-0.5}, {#2}) rectangle ({#1+0.5}, {#2-1.0});
    \draw ({#1-0.5}, {#2-1.0}) -- ({#1+0.5}, {#2});
    % Join line that crosses the indicator (mid-point y-0.5)
    \draw ({#1+0.5}, {#2-0.5}) -- ({#1+1.2}, {#2-0.5});
    % Hourglass valve indicator (two triangles meeting at center)
    \draw ({#1+0.8}, {#2-0.25}) -- ({#1+1.2}, {#2-0.25}) -- ({#1+0.8}, {#2-0.75}) -- ({#1+1.2}, {#2-0.75}) -- cycle;
    % Label
    \node[anchor=south] at ({#1+1.0}, {#2-0.25}) {#3};
}
